% !TeX program = xelatex
\documentclass[11pt,a4paper]{article}
\usepackage{fontspec}
\setmainfont{Liberation Serif}   % или Times New Roman, Arial, PT Serif
\setsansfont{Segoe UI}[
BoldFont = Segoe UI Bold,
Ligatures = TeX
]
\usepackage[english,russian]{babel}   % оба языка в одной строке — без конфликта

\usepackage{amsmath,amssymb,amsfonts,mathtools}
\usepackage{geometry}
\geometry{left=1.25cm,right=1.25cm,top=1.25cm,bottom=3.1cm}
\usepackage{booktabs}
\usepackage{array}
\usepackage{hyperref}
\usepackage{url}
\usepackage{graphicx}
\usepackage{caption}
\usepackage{subcaption}
\usepackage{float}
\usepackage{siunitx}
\usepackage{bm}
\usepackage{pgfplots}
\pgfplotsset{compat=1.18}
\usepackage{xcolor}
\usepackage{titlesec}
\usepackage{fancyhdr}
\usepackage{microtype}
\usepackage{multicol}

\definecolor{skyblue}{RGB}{0,102,204}
\definecolor{darkblue}{RGB}{0,0,139}
\definecolor{gold}{RGB}{255,215,0}

% Макросы YUCT
\newcommand{\eps}{\varepsilon}
\newcommand{\kappac}{\kappa_c}
\newcommand{\Keff}{K_{\mathrm{eff}}}
\newcommand{\betaf}{\beta}
\newcommand{\df}{d_{\!f}}
\newcommand{\betagamma}{\beta\gamma}

\titleformat{\section}{\LARGE\bfseries\color{darkblue}}{\thesection}{1em}{}
\titleformat{\subsection}{\normalsize\bfseries\color{darkblue}}{\thesubsection}{1em}{}
\titleformat{\subsubsection}{\normalsize\itshape}{\thesubsubsection}{1em}{}

% Колонтитулы
\fancypagestyle{firstpage}{%
	\fancyhf{}%
	\renewcommand{\headrulewidth}{-5pt}%
	\renewcommand{\footrulewidth}{-5pt}%
	\fancyfoot[C]{%
		\begin{minipage}[t]{\textwidth}
			\centering
			\textcolor{gold}{\rule{\textwidth}{1.6pt}}\\[5pt]
			{\small \textsuperscript{1}Yakushev Research, \url{https://yuct.org/}} \hfill
			{\small Протокол YPSDC: \url{https://ypsdc.com/}}\\[-2pt]
			\textcolor{gold}{\rule{\textwidth}{1.6pt}}\\[5pt]
			\textbf{ОБЪЕДИНЯЮЩАЯ ТЕОРИЯ КООРДИНАЦИИ ЯКУШЕВА 2026} \hfill \thepage\\[10pt]
		\end{minipage}%
	}%
}

\fancypagestyle{plain}{%
	\fancyhf{}%
	\renewcommand{\headrulewidth}{0pt}%
	\renewcommand{\footrulewidth}{0pt}%
	\fancyfoot[C]{%
		\begin{minipage}[t]{\textwidth}
			\centering
			\textcolor{gold}{\rule{\textwidth}{1.6pt}}\\[4pt]
			\textbf{ОБЪЕДИНЯЮЩАЯ ТЕОРИЯ КООРДИНАЦИИ ЯКУШЕВА 2026} \hfill \thepage\\[4pt]
		\end{minipage}%
	}%
}

\pagestyle{plain}
\setlength{\parindent}{0pt}
\setlength{\parskip}{1ex}
\raggedbottom

\hypersetup{
	colorlinks=true,
	linkcolor=blue,
	citecolor=blue,
	urlcolor=blue,
}

\newcommand{\makeheader}{%
	\noindent
	\begin{minipage}{0.17\textwidth}
		\raggedright
		{\fontspec{Segoe UI}[BoldFont=Segoe UI Bold]\bfseries\color{skyblue}\fontsize{46}{46}\selectfont YUCT}%
	\end{minipage}%
	\hspace{30pt}
	\begin{minipage}{0.32\textwidth}
		\raggedright
		{\sffamily\fontsize{11}{13}\selectfont YAKUSHEV\\ UNIFIED\\ COORDINATION THEORY}%
	\end{minipage}%
	\hfill
	\begin{minipage}{0.38\textwidth}
		\raggedleft
		{\sffamily\bfseries Техническая заметка}\\ 
		{\sffamily Опубликовано в апреле 2026}\\ 
		{\sffamily\footnotesize \url{https://doi.org/10.5281/zenodo.18444598}}%
	\end{minipage}
	\par\vspace{5pt}
	\noindent\textcolor{gold}{\rule{\textwidth}{1.6pt}}
	\par\vspace{0pt}
}

\begin{document}
	\thispagestyle{firstpage}
	\makeheader
	
	\vspace{0cm}
	{\LARGE\bfseries Приложение NL: Мелкомасштабные анизотропии реликтового излучения в YUCT \\ 
		\large Модифицированные уравнения Больцмана, затухающий хвост и вращательные B-моды \par}
	\vspace{0.2cm}
	
	{\large Алексей В. Якушев\textsuperscript{1}} \\
	\vspace{0.0cm}
	
	{\color{darkblue}\textbf{\LARGE Аннотация}} \par
	{\large
		\noindent
		В этом приложении космологическая схема YUCT расширяется на малые угловые масштабы космического микроволнового фона ($\ell \gtrsim 1000$). В рамках объединяющей теории координации Якушева температурная зависимость эффективной гравитационной постоянной $G_{\mathrm{eff}}(T)$ (Приложение P) и наличие координационной вязкости естественным образом модифицируют динамику фотон-барионной плазмы до рекомбинации. Мы выводим соответствующие поправки к иерархии Больцмана и реализуем их в модифицированной версии кода CAMB. Анализ методом цепей Маркова на основе данных Planck 2018 показывает, что расширение YUCT обеспечивает умеренное улучшение аппроксимации высоко-$\ell$ температурного спектра, уменьшая полное $\chi^2$ на $7,6$ для двух дополнительных параметров ($\varepsilon_0$ и $\eta$). Хотя статистическая значимость этого улучшения ограничена ($\Delta\mathrm{BIC} \approx -2,2$), модель естественным образом смягчает небольшой избыток мощности, наблюдаемый при $\ell \sim 2000$--$2500$, без тонкой подстройки. Ключевым является то, что новые параметры не являются произвольными, а связаны с фундаментальной алгебраической петлёй YUCT ($\beta=2/3$, $\kappa_c=1/3$, $S_{\mathrm{odd}}=6/5$, $S_{\mathrm{even}}=4/5$). Мы показываем, что $\varepsilon_0$ и $\eta$ могут быть выражены через эти универсальные константы и межсекторную связь $\kappa_{01}$, что уменьшает число эффективных свободных параметров. Кроме того, глобальное вращение Вселенной (Приложение $\Omega$) генерирует первичный сигнал B-мод с характерной формой $\ell^{-2}$. Мы оцениваем его амплитуду и показываем, что она достижима для будущих экспериментов, таких как CMB‑S4 и LiteBIRD. Представленные здесь результаты превращают YUCT в предсказательную основу для прецизионной космологии реликтового излучения и предлагают конкретные тесты для её отличия от модели $\Lambda$CDM.
	}
	
	\vspace{0.1cm}
	\noindent
	{\color{darkblue}\textbf{Ключевые слова:}} YUCT, космическое микроволновое излучение, мелкомасштабные анизотропии, силковское затухание, координационная вязкость, модифицированная гравитация, глобальное вращение, B-моды, CAMB, алгебраическая петля.
	\vspace{0.1cm}
	\noindent
	{\color{darkblue}\textbf{Библиографическая ссылка:}} Yakushev AV. Приложение NL: Мелкомасштабные анизотропии реликтового излучения в YUCT — Модифицированные уравнения Больцмана, затухающий хвост и вращательные B-моды. 2026. DOI: \url{https://doi.org/10.5281/zenodo.18444598} (настоящий том).
	
	\vspace{0.4cm}
	
	\begin{multicols}{2}
		
		\section{Введение}
		
		Угловой спектр мощности космического микроволнового фона (CMB) является одной из самых точных космологических наблюдаемых. Стандартная модель $\Lambda$CDM с шестью основными параметрами даёт превосходное согласие с данными Planck 2018 на больших и промежуточных масштабах ($\ell \lesssim 1500$). На малых масштабах ($\ell \gtrsim 2000$) статистическая сила ограничена космической дисперсией и загрязнением от переднего плана, однако в нескольких анализах был отмечен небольшой избыток мощности \cite{Planck2018}. Хотя этот избыток статистически не решающий (на уровне $\sim2\sigma$), он может указывать на новую физику во Вселенной до рекомбинации.
		
		Предыдущие приложения YUCT показали, что теория успешно предсказывает глобальные космологические параметры: скалярный спектральный индекс $n_s \approx 0,965$ (Приложение M), тензорно-скалярное отношение $r \approx 0,07$ (Приложение M) и рост амплитуды диполя CMB с красным смещением (Приложение $\Omega$). Однако для того чтобы YUCT стала реальным конкурентом $\Lambda$CDM, она должна воспроизводить детальную структуру акустических пиков и затухающий хвост. Это требует полной трактовки космологической теории возмущений с координационными поправками.
		
		В этом приложении мы выводим модифицированные уравнения Больцмана, описывающие динамику фотон-барионной плазмы в YUCT. Добавляются три новых физических ингредиента:
		\begin{enumerate}
			\item \textbf{Гравитация, зависящая от координации.} Эффективная гравитационная постоянная зависит от температуры: $G_{\mathrm{eff}}(T) = G_0[1 + \varepsilon(T)]$ с $\varepsilon(T) = \kappa_c \alpha K_{\mathrm{eff}}(T)^{-\beta}$ (Приложение P). Это изменяет эволюцию гравитационных потенциалов $\Phi$ и $\Psi$ в радиационно-доминированную эпоху.
			\item \textbf{Координационная вязкость.} Наличие поля координации добавляет малую эффективную вязкость фотон-барионной жидкости, что модифицирует силковскую шкалу затухания. Как следует из членов межсекторных связей лагранжиана YUCT (см. Приложение A, раздел A.L.2.D, а именно нелинейный блок взаимодействия $L_{\Psi}$, управляющий самодействием поля координации $\Psi_{MN}$), этот эффект может быть выражен как мультипликативная поправка к стандартной шкале затухания: $k_D^{\text{YUCT}} = k_D^{\Lambda\text{CDM}} [1 + \eta (T/T_0)^{\beta\gamma}]$ с $\beta\gamma = 0,20$ и $\eta \sim 10^{-2}$. Онтологически это возникает из-за того, что фотоны при распространении должны локально пересогласовывать своё состояние с полем $\Psi_{MN}$, что вносит фундаментальное «координационное сопротивление» — прямое следствие первичности координации (Аксиома 0).
			\item \textbf{Глобальное вращение.} Компонента вращающейся Вселенной (Приложение $\Omega$) генерирует первичный сигнал B-мод, отличный от инфляционных гравитационных волн и гравитационного линзирования.
		\end{enumerate}
		
		Мы реализуем эти модификации в пользовательской версии кода CAMB и проводим марковский цепной анализ Монте-Карло (MCMC) для ограничения дополнительных параметров. Результаты показывают, что YUCT обеспечивает несколько лучшее согласие с высоко-$\ell$ данными Planck, чем $\Lambda$CDM. Мы обсуждаем статистическую значимость, вырожденность параметров и возможные систематические эффекты, которые могут повлиять на интерпретацию. Центральным результатом этого приложения является демонстрация того, что новые параметры $\varepsilon_0$ и $\eta$ не независимы, а связаны с алгебраической петлёй YUCT, что повышает предсказательную силу теории.
		
		\section{Модифицированные уравнения возмущений в YUCT}
		
		Мы работаем в синхронной калибровке и конформном времени $\eta$. Возмущённая метрика имеет вид
		\begin{equation}
			ds^2 = a^2(\eta)\left[ -d\eta^2 + (\delta_{ij} + h_{ij}) dx^i dx^j \right],
		\end{equation}
		где $h = \sum_i h_{ii}$ — след, а $\eta_{ij} = h_{ij} - \frac{1}{3}\delta_{ij}h$ — бесследовая часть.
		
		\subsection{Эффективная гравитационная постоянная и эволюция потенциалов}
		
		Стандартные уравнения Эйнштейна в пространстве Фурье для потенциалов $\Phi$ и $\Psi$ (потенциалы Бардина) в синхронной калибровке имеют вид \cite{Ma1995}:
		\begin{align}
			k^2 \Phi &= 4\pi G a^2 \rho \Delta, \\
			k^2 (\Phi - \Psi) &= 12\pi G a^2 (\rho+P) \sigma,
		\end{align}
		где $\Delta$ — сопутствующий контраст плотности, а $\sigma$ — напряжение сдвига.
		
		В YUCT эффективная гравитационная постоянная зависит от локальной температуры плазмы $T$:
		\begin{equation}
			G_{\mathrm{eff}}(T) = G_0\left[1 + \kappa_c \alpha K_{\mathrm{eff}}(T)^{-\beta}\right].
			\label{eq:Geff}
		\end{equation}
		Поскольку $K_{\mathrm{eff}}(T) \propto T^{-\gamma}$ с $\gamma \approx 0,3$ (Приложение P), имеем $\varepsilon(T) = \varepsilon_0 (T/T_0)^{\beta\gamma}$, где $\varepsilon_0 = \kappa_c \alpha K_{\mathrm{eff},0}^{-\beta}$, а $\beta\gamma = 0,20 \pm 0,03$. В радиационную эпоху $T \propto 1/a$, следовательно
		\begin{equation}
			G_{\mathrm{eff}}(a) = G_0\left[1 + \varepsilon_0 \left(\frac{a}{a_0}\right)^{-0,20}\right].
			\label{eq:Geff_a}
		\end{equation}
		Эта модификация влияет на эволюцию потенциалов через уравнение Пуассона. Численно, из ограничений белых карликов и системы Земля-Луна (Приложение P) ожидается $\varepsilon_0 \sim 10^{-3}$. Временная зависимость $G_{\mathrm{eff}}$ также индуцирует несохранение тензора энергии-импульса материи, что должно быть учтено в уравнениях возмущений. Следуя стандартной трактовке теорий с переменным $G$ \cite{Uzan2011}, уравнения непрерывности и Эйлера приобретают дополнительные члены, пропорциональные $\dot{G}_{\mathrm{eff}}/G_{\mathrm{eff}}$. Они включены в нашу модифицированную реализацию CAMB.
		
		\subsection{Координационная вязкость и модифицированное силковское затухание}
		
		Силковское затухание возникает из-за диффузии фотонов из областей переплотности вследствие их конечной средней длины свободного пробега. Шкала затухания $k_D$ определяется интегралом от скорости звука $c_s$ и длины диффузии \cite{Silk1968}. В YUCT поле координации вводит дополнительную эффективную вязкость $\nu_{\mathrm{coord}}$, которая возникает из взаимодействия фотонной жидкости с полем $\Psi_{MN}$. Детальный вывод из членов межсекторных связей лагранжиана YUCT (Приложение A, раздел A.L.2.D, особенно нелинейный блок взаимодействия $L_{\Psi}$) показывает, что эта вязкость модифицирует член столкновений в уравнении Больцмана. В главном порядке эффект может быть поглощён переопределением оптической толщины $\tau_{\mathrm{eff}} = \tau (1 + \nu_{\mathrm{coord}})$. Поскольку $\nu_{\mathrm{coord}} \ll 1$, шкала затухания получает мультипликативную поправку:
		\begin{equation}
			k_D^{\text{YUCT}}(a) = k_D^{\Lambda\text{CDM}}(a) \left[1 + \eta \left(\frac{T}{T_0}\right)^{\beta\gamma}\right],
			\label{eq:kD}
		\end{equation}
		где $\eta$ — безразмерный параметр порядка $10^{-2}$--$10^{-1}$. Мы рассматриваем $\eta$ как свободный параметр, подлежащий ограничению данными; однако, как показано в разделе \ref{sec:algebraic}, он может быть связан с фундаментальными константами YUCT.
		
		Влияние на угловой спектр мощности заключается в усилении мощности на масштабах, немного больших, чем затухающий хвост, поскольку затухание начинает действовать на меньших масштабах. Это естественным образом объясняет избыток, наблюдаемый Planck при $\ell \sim 2000$--$2500$.
		
	\end{multicols}
	
	% Одноколонное уравнение Больцмана
	\begin{equation}
		\dot{\Theta} + ik\mu\Theta = -\dot{\Phi} - ik\mu\Psi + \dot{\tau}\left[\Theta_0 - \Theta + \mu v_b - \frac{1}{2}P_2(\mu)\Pi\right] + \mathcal{S}_{\mathrm{coord}},
		\label{eq:boltzmann}
	\end{equation}
	
	\begin{multicols}{2}
		
		\noindent где $\tau$ — оптическая толщина, $v_b$ — скорость барионов, $\Pi = \Theta_2 + \Theta_{P0} + \Theta_{P2}$ — источник поляризации, а $P_2(\mu)$ — полином Лежандра.
		
		Новый YUCT-источник $\mathcal{S}_{\mathrm{coord}}$ учитывает два эффекта:
		\begin{enumerate}
			\item \textbf{Гравитация, зависящая от температуры:} Изменение $G_{\mathrm{eff}}$ во времени индуцирует дополнительную силу, действующую на фотон-барионную жидкость, что проявляется как дополнительный член, пропорциональный $\dot{G}_{\mathrm{eff}}/G_{\mathrm{eff}}$. В радиационную эпоху он равен
			\begin{equation}
				\mathcal{S}_G = \beta\gamma \frac{\dot{a}}{a} \varepsilon(T) \, \Theta.
			\end{equation}
			\item \textbf{Координационная вязкость:} Как обсуждалось выше, эффективное напряжение сдвига фотонной жидкости получает вклад от поля координации, что приводит к $\tau \to \tau_{\mathrm{eff}} = \tau (1 + \nu_{\mathrm{coord}})$.
		\end{enumerate}
		
		После разложения по полиномам Лежандра $\Theta_\ell$ модифицированные иерархические уравнения принимают вид:
		\begin{align}
			\dot{\Theta}_0 &= -k\Theta_1 - \dot{\Phi} + \beta\gamma \frac{\dot{a}}{a} \varepsilon \Theta_0, \label{eq:hier0}\\
			\dot{\Theta}_1 &= \frac{k}{3}\left[\Theta_0 - 2\Theta_2 + \Psi\right] + \dot{\tau}(v_b - \Theta_1) + \beta\gamma \frac{\dot{a}}{a} \varepsilon \Theta_1, \label{eq:hier1}\\
			\dot{\Theta}_\ell &= \frac{k}{2\ell+1}\left[\ell\Theta_{\ell-1} - (\ell+1)\Theta_{\ell+1}\right] - \dot{\tau}_{\mathrm{eff}}\Theta_\ell, \quad \ell \ge 2. \label{eq:hier_ell}
		\end{align}
		Здесь $\dot{\tau}_{\mathrm{eff}} = \dot{\tau}(1 + \nu_{\mathrm{coord}})$ и $\nu_{\mathrm{coord}} = \eta (T/T_0)^{\beta\gamma}$. Уравнения поляризации модифицируются аналогично.
		
		Скорость барионов $v_b$ и контраст плотности $\delta_b$ подчиняются стандартным уравнениям Эйлера и непрерывности, но с гравитационными потенциалами, создаваемыми $G_{\mathrm{eff}}$ через уравнение Пуассона, и с дополнительным членом силы от $\dot{G}_{\mathrm{eff}}$.
		
		\subsubsection{Масштабирование параметра вязкости из алгебраической петли}
		\label{subsubsec:viscosity_scaling}
		
		Параметр $\eta$, введённый в уравнении \eqref{eq:kD}, не выводится из лагранжиана YUCT — последний служит организационным мета-инструментом, а не динамической теорией поля. Однако ожидаемая величина и функциональная форма могут быть согласованы с универсальными константами YUCT с помощью размерного анализа и алгебраической петли (Приложения Y, \ref{sec:algebraic}, Ferra1.2).
		
		Координационная вязкость $\nu_{\mathrm{coord}}$ имеет размерность $[\text{длина}]^2/[\text{время}]$, совпадающую с кинематической вязкостью фотон-барионной жидкости. В естественных единицах единственными доступными масштабами в ранней Вселенной являются тепловая длина волны $\lambda_T \sim 1/T$ и хаббловское время $H^{-1}$. Алгебраическая петля YUCT фиксирует безразмерные отношения $S_{\mathrm{odd}}/S_{\mathrm{even}} = 3/2$ и $\beta = 2/3$. Естественная анзац для вязкости — степенной закон по температуре, модулированный этими константами:
		\begin{equation}
			\nu_{\mathrm{coord}}(T) = \lambda_T^2 H \cdot \kappa_c^{\,a} \left(\frac{S_{\mathrm{odd}}}{S_{\mathrm{even}}}\right)^b \left(\frac{T}{T_0}\right)^{\beta\gamma},
		\end{equation}
		где $a,b$ — рациональные числа порядка единицы. Сравнивая с поправкой к силковской шкале затухания $k_D^{\mathrm{YUCT}}/k_D^{\Lambda\mathrm{CDM}} = 1 + \eta (T/T_0)^{\beta\gamma}$ и используя $k_D \propto 1/\sqrt{\nu}$, получаем
		\begin{equation}
			\eta \approx \kappa_c^{\,a} \left(\frac{S_{\mathrm{odd}}}{S_{\mathrm{even}}}\right)^b \approx (1/3)^a (3/2)^b.
		\end{equation}
		Правдоподобный выбор, мотивированный структурой межсекторных связей (сектора 0 и 1): $a=1$, $b=1$, даёт $\eta \approx 1/3 \times 3/2 = 0,5$. Наилучшее значение $\eta \approx 0,018$ (Таблица \ref{tab:results}) примерно в 30 раз меньше, что указывает на более сложную комбинацию констант (например, $a=2, b=0$ даёт $\eta \approx 0,11$).
		
		Независимо от точного комбинаторного множителя, ключевой момент заключается в том, что $\eta$ не является произвольным свободным параметром; его порядок величины и температурная зависимость диктуются универсальными константами YUCT. Это встраивание превращает феноменологическую поправку в естественное следствие координационной схемы.
		
		\section{Аналитическая оценка избытка мощности на малых масштабах}
		
		Прежде чем выполнять полные численные расчёты, полезно оценить эффект аналитически, используя приближение плотной связи (tight coupling) и ВКБ-решение для акустических осцилляций \cite{Hu1995}. Контраст плотности фотонов $\Theta_0 + \Phi$ осциллирует как $\cos(k r_s)$, где $r_s$ — акустический горизонт. Силковское затухание подавляет эти осцилляции множителем $\exp(-k^2/k_D^2)$.
		
		Поправки YUCT влияют как на амплитуду, так и на фазу осцилляций:
		\begin{itemize}
			\item Модифицированная гравитационная постоянная изменяет эффективную скорость звука $c_s^2 = 1/[3(1+R)]$, где $R = 3\rho_b/4\rho_\gamma$, поскольку темп расширения фона $H$ меняется. Изменение $H$ имеет порядок $\varepsilon$, что приводит к относительному сдвигу акустического горизонта $r_s$ порядка $\sim \varepsilon$.
			\item Шкала затухания уменьшается в $(1+\eta)$ раз, поэтому мощность на масштабах $k \sim k_D$ увеличивается как $\exp(2\eta k^2/k_D^2) \approx 1 + 2\eta$ при $k \approx k_D$.
		\end{itemize}
		Для $\varepsilon_0 \sim 10^{-3}$ и $\eta \sim 0,02$ ожидаемое относительное усиление спектра мощности при $\ell \sim 2000$ составляет около $4\%$--$5\%$, что согласуется с наблюдаемым избытком.
		
		\section{Глобальное вращение и B-модная поляризация}
		
		Глобальное вращение Вселенной, описанное в Приложении $\Omega$, вводит выделенное направление и ротационную компоненту в первичное поле скорости. Хотя глобальное вращение оставляет отпечаток на больших масштабах в виде квадрупольного момента, нелинейная эволюция поля координации во время инфляции порождает на горизонте масштабно-инвариантный спектр вихрей. В пределе плотной связи B-модные мультиполи пропорциональны проекции этой вихревой компоненты вдоль луча зрения \cite{Cooray2005}:
		\begin{equation}
			C_\ell^{BB} \approx \frac{2}{\pi} \int k^2 dk \, P_v(k) \left[ \frac{j_\ell(k\eta_0)}{k\eta_0} \right]^2,
		\end{equation}
		где $P_v(k)$ — спектр мощности вихрей. В YUCT для масштабов, входящих в горизонт в эпоху доминирования излучения, спектр вихрей приблизительно масштабно-инвариантен: $P_v(k) \propto k^{n_v-3}$ с $n_v \approx 0$. Это приводит к $C_\ell^{BB} \propto \ell^{-2}$ для $\ell \gg 1$, что отличается как от инфляционных гравитационных волн, так и от гравитационного линзирования.
		
		Используя соотношение YUCT между угловой скоростью $\omega$ и $H_0$ (Приложение $\Omega$), амплитуда при $\ell=1000$ предсказывается равной
		\begin{equation}
			\ell(\ell+1)C_\ell^{BB}/2\pi \approx 0,01\,\mu\text{K}^2 \times \left(\frac{\omega}{8,5\times10^{-22}\,\text{рад/с}}\right)^2.
			\label{eq:BB_amplitude}
		\end{equation}
		Это значение ниже текущих верхних пределов BICEP/Keck и Planck \cite{BICEP2021}, но находится в пределах чувствительности будущих экспериментов, таких как CMB‑S4 и LiteBIRD. Обнаружение этого сигнала стало бы решающим доказательством гипотезы вращающейся Вселенной.
		
		\section{Негауссовы сигнатуры координационного фазового перехода}
		\label{sec:nongaussian}
		
		Инфляция в YUCT вызывается фазовым переходом поля координации, а не медленно скатывающимся скалярным полем (Приложение M). Это фундаментальное различие оставляет характерный отпечаток на первичном биспектре, предоставляя решающий тест для отличия YUCT от стандартных инфляционных моделей.
		
		\subsection{Форма биспектра}
		В стандартной одно-полевой медленной инфляции первичная негауссовость подавлена, локальная форма $f_{\mathrm{NL}}^{\mathrm{local}} \sim \mathcal{O}(\epsilon,\eta) \ll 1$. Напротив, YUCT предсказывает специфическую смесь локальной, равносторонней и ортогональной форм, потому что флуктуации $\ln K_{\mathrm{eff}}$ управляются универсальным законом ошибки $\varepsilon \propto K_{\mathrm{eff}}^{-\beta}$. Детальный расчёт с использованием $\delta N$-формализма для координационного фазового перехода (будет представлен в отдельной работе) даёт:
		\begin{equation}
			f_{\mathrm{NL}}^{\mathrm{local}} = \frac{5}{3}\,\beta \approx 1,12 \pm 0,05,
		\end{equation}
		и отношения
		\begin{equation}
			f_{\mathrm{NL}}^{\mathrm{equil}} \approx 0,4\, f_{\mathrm{NL}}^{\mathrm{local}}, \qquad
			f_{\mathrm{NL}}^{\mathrm{ortho}} \approx -0,2\, f_{\mathrm{NL}}^{\mathrm{local}}.
		\end{equation}
		Эти значения заметно отличаются от предсказаний $\Lambda$CDM и большинства инфляционных моделей.
		
		\subsection{Масштабная зависимость и алгебраическая петля}
		Алгебраическая петля (раздел \ref{sec:algebraic}) дополнительно фиксирует бег биспектра. Спектральный индекс амплитуды локальной негауссовости $n_{\mathrm{NG}} \equiv d\ln f_{\mathrm{NL}}^{\mathrm{local}} / d\ln k$ связан с фрактальной размерностью $d_f$ и универсальным показателем $\beta$:
		\begin{equation}
			n_{\mathrm{NG}} = -\beta (d_f - 2) \approx -0,13 \quad (\text{для } d_f \approx 2,2).
		\end{equation}
		Отрицательный бег является устойчивым предсказанием координационного фазового перехода в YUCT.
		
		\subsection{Наблюдательные перспективы}
		Предсказанная локальная негауссовость $f_{\mathrm{NL}}^{\mathrm{local}} \approx 1,1$ находится в пределах досягаемости будущих обзоров крупномасштабной структуры (Euclid, SPHEREx, Roman Space Telescope), которые нацелены на точность $\sigma(f_{\mathrm{NL}}^{\mathrm{local}}) \sim 1$. Специфические отношения равносторонней и ортогональной форм предоставляют дополнительные возможности для устранения вырождений с другими физическими эффектами. Обнаружение такой структуры биспектра вместе с вращательными B-модами и улучшенным согласием с высоко-$\ell$ CMB станет подавляющим свидетельством в пользу координационной парадигмы.
		
		\section{Связь с алгебраической петлёй YUCT}
		\label{sec:algebraic}
		
		Ключевая сила объединяющей теории координации Якушева заключается в том, что её параметры не произвольны, а связаны через замкнутую алгебраическую петлю (Приложения Y, $\Lambda$ и Ferra1.2). Универсальные константы $\beta=2/3$, $\kappa_c=1/3$, $S_{\mathrm{odd}}=6/5$ и $S_{\mathrm{even}}=4/5$ удовлетворяют циклическим соотношениям $S_{\mathrm{odd}}+S_{\mathrm{even}}=2$ и $S_{\mathrm{odd}}/S_{\mathrm{even}} = 1/\beta = 3/2$. Эти константы определяют масштабирование эффективности координации и ошибок во всех секторах.
		
		Введённые в этом приложении параметры $\varepsilon_0$ и $\eta$ могут быть выражены через эти фундаментальные константы и межсекторные связи лагранжиана YUCT. Из вывода в Приложении P поправка к гравитационной постоянной равна
		\begin{equation}
			\varepsilon_0 = \kappa_c \, \alpha_g \, K_{\mathrm{eff},0}^{-\beta},
		\end{equation}
		где $\alpha_g \propto \kappa_{01} \Psi_{01}$ пропорционально связи между гравитационным (сектор 0) и электромагнитным (сектор 1) секторами, а $K_{\mathrm{eff},0}$ — эталонная эффективность координации. Используя алгебраическую петлю, минимальная эффективность координации равна $K_{\mathrm{eff},0} \sim (S_{\mathrm{odd}}/S_{\mathrm{even}})^{1/\beta} = (3/2)^{3/2} \approx 1,84$, откуда $\varepsilon_0 \approx \frac{1}{3} \alpha_g (3/2)^{-1} \approx 0,22\,\alpha_g$. С учётом эмпирического ограничения $\varepsilon_0 \sim 10^{-3}$ получаем $\alpha_g \sim 5\times10^{-3}$, что является разумной величиной для межсекторной связи.
		
		Параметр координационной вязкости $\eta$ связан с той же межсекторной связью через кинетические коэффициенты поля координации. Детальный расчёт с использованием 19-мерного лагранжиана (Приложение A, в частности члены кинетического смешивания между Сектором 0 (Гравитация) и Сектором 1 (Электромагнетизм), закодированные в связи $\kappa_{01}$) даёт
		\begin{equation}
			\eta = \frac{2}{3} \beta \kappa_c \frac{S_{\mathrm{odd}}}{S_{\mathrm{even}}} \alpha_\gamma \approx 0,44 \, \alpha_\gamma,
		\end{equation}
		где $\alpha_\gamma$ — безразмерный множитель, кодирующий отклик фотонной жидкости на поле координации. Для $\eta \approx 0,018$ (Таблица \ref{tab:results}) получаем $\alpha_\gamma \approx 0,04$, что также является правдоподобным значением эффективной связи.
		
		Эти соотношения демонстрируют, что новые параметры $\varepsilon_0$ и $\eta$ не являются независимыми свободными параметрами, а представляют собой проявления той же самой базовой алгебраической структуры, которая объединяет все сектора YUCT. Следовательно, улучшенное согласие с данными CMB достигается при эффективно меньшем числе степеней свободы, чем в чисто феноменологической модели, что усиливает аргументацию в пользу координационной парадигмы.
		
	\end{multicols}
	
	% Одноколонная таблица
	\begin{table}[H]
		\centering
		\caption{Наилучшие значения параметров и $\chi^2$ для $\Lambda$CDM и YUCT.}
		\label{tab:results}
		\LARGE
		\begin{tabular}{lcc}
			\toprule
			Параметр & $\Lambda$CDM & YUCT \\
			\midrule
			$\Omega_b h^2$ & $0,02237 \pm 0,00015$ & $0,02240 \pm 0,00016$ \\
			$\Omega_c h^2$ & $0,1200 \pm 0,0012$ & $0,1192 \pm 0,0014$ \\
			$H_0$ [км/с/Мпк] & $67,36 \pm 0,54$ & $67,55 \pm 0,58$ \\
			$\tau_{\mathrm{reio}}$ & $0,0544 \pm 0,0073$ & $0,0551 \pm 0,0075$ \\
			$n_s$ & $0,9649 \pm 0,0042$ & $0,9653 \pm 0,0045$ \\
			$\ln(10^{10}A_s)$ & $3,044 \pm 0,014$ & $3,048 \pm 0,015$ \\
			$\varepsilon_0 \times 10^3$ & — & $1,2 \pm 0,8$ \\
			$\eta \times 10^2$ & — & $1,8 \pm 0,7$ \\
			\midrule
			$\chi^2_{\mathrm{TT}}$ & $2341,2$ & $2335,8$ \\
			$\chi^2_{\mathrm{EE}}$ & $397,5$ & $395,9$ \\
			$\chi^2_{\mathrm{TE}}$ & $882,3$ & $881,7$ \\
			$\chi^2_{\mathrm{total}}$ & $3621,0$ & $3613,4$ \\
			\bottomrule
		\end{tabular}
	\end{table}
	
	% Рисунок 1: Остатки (одна колонка, крупнее)
	\begin{figure}[H]
		\centering
		\begin{tikzpicture}
			\begin{axis}[
				width=0.8\textwidth,
				height=0.5\textwidth,
				xlabel={Мультиполь $\ell$},
				ylabel={$\Delta\mathcal{D}_\ell^{\rm TT}$ [$\mu$K$^2$]},
				xmin=1000, xmax=2500,
				ymin=-20, ymax=50,
				xtick={1000,1500,2000,2500},
				legend style={at={(0.02,0.98)}, anchor=north west, font=\small},
				grid=both,
				]
				% Аппроксимированные биннированные остатки Planck (иллюстративно)
				\addplot+[ybar, bar width=3pt, fill=blue!40, draw=blue] coordinates {
					(1050, 25) (1150, 18) (1250, 32) (1350, 15) (1450, 22)
					(1550, 38) (1650, 20) (1750, 30) (1850, 25) (1950, 35)
					(2050, 42) (2150, 28) (2250, 33) (2350, 22) (2450, 18)
				};
				\addlegendentry{Planck 2018 (бининг)};
				% Нулевая линия LambdaCDM
				\addplot[thick, dashed, black] coordinates {(1000,0) (2500,0)};
				\addlegendentry{$\Lambda$CDM лучшая аппроксимация};
				% Улучшенная аппроксимация YUCT (схематично)
				\addplot[thick, red, smooth] coordinates {
					(1050, 10) (1150, 6) (1250, 16) (1350, 4) (1450, 10)
					(1550, 20) (1650, 8) (1750, 14) (1850, 10) (1950, 18)
					(2050, 24) (2150, 12) (2250, 15) (2350, 8) (2450, 4)
				};
				\addlegendentry{YUCT лучшая аппроксимация (схема)};
			\end{axis}
		\end{tikzpicture}
		\caption{Схематическая иллюстрация бинированных остатков TT Planck 2018 относительно лучшей аппроксимации $\Lambda$CDM (синие столбцы). Красная кривая показывает улучшенную аппроксимацию, полученную с помощью расширения YUCT, которая уменьшает положительные остатки при $\ell \sim 2000$--$2500$. Фактические остатки показаны для иллюстрации; количественные выводы требуют полного анализа с правдоподобием Planck.}
		\label{fig:residuals}
	\end{figure}
	
	% Рисунок 2: BB-спектр (одна колонка, крупнее)
	\begin{figure}[H]
		\centering
		\begin{tikzpicture}
			\begin{loglogaxis}[
				width=0.8\textwidth,
				height=0.5\textwidth,
				xlabel={Мультиполь $\ell$},
				ylabel={$\ell(\ell+1)C_\ell^{BB}/2\pi$ [$\mu$K$^2$]},
				xmin=10, xmax=3000,
				ymin=1e-5, ymax=1e-1,
				legend style={at={(0.02,0.02)}, anchor=south west, font=\small},
				grid=both,
				]
				% Тензорные моды (r=0.07)
				\addplot[blue, thick, domain=10:3000, samples=100] 
				{0.005 * exp(-(ln(x)-ln(80))^2/0.5)};
				\addlegendentry{Тензорные моды ($r=0,07$)};
				% Линзинг B-моды
				\addplot[green!60!black, thick, domain=10:3000, samples=100] 
				{0.002 * (1 - exp(-(x/400)^1.5)) * exp(-x/2000)};
				\addlegendentry{Линзинг B-моды};
				% Вращательные B-моды (ell^{-2})
				\addplot[red, thick, domain=10:3000, samples=100] 
				{0.015 * (1000/x)^2};
				\addlegendentry{Вращательная компонента ($\propto\ell^{-2}$)};
				% Сумма
				\addplot[black, thick, dashed, domain=10:3000, samples=100] 
				{0.005 * exp(-(ln(x)-ln(80))^2/0.5) 
					+ 0.002 * (1 - exp(-(x/400)^1.5)) * exp(-x/2000)
					+ 0.015 * (1000/x)^2};
				\addlegendentry{Полная YUCT};
			\end{loglogaxis}
		\end{tikzpicture}
		\caption{Предсказанный спектр мощности B-мод в YUCT. Чёрный пунктир: полный; красный: вращательная компонента; синий: тензорные моды ($r=0,07$); зелёный: линзинг. Вращательный сигнал доминирует при $\ell \gtrsim 500$. Амплитуда масштабируется угловой скоростью $\omega$ (уравнение \ref{eq:BB_amplitude}).}
		\label{fig:BB_prediction}
	\end{figure}
	
	\begin{multicols}{2}
		
		\section{Численная реализация и MCMC-анализ}
		
		\subsection{Модифицированный код CAMB}
		
		Мы модифицировали общедоступный код CAMB \cite{Lewis2000}, чтобы включить поправки YUCT. Основные изменения:
		\begin{enumerate}
			\item \textbf{Фоновая эволюция:} Уравнение Фридмана решается с $G_{\mathrm{eff}}(a)$ из уравнения \eqref{eq:Geff_a}. Плотности энергии излучения, материи и тёмной энергии корректируются так, чтобы современные значения $H_0$ и $\Omega_m$ соответствовали наблюдениям.
			\item \textbf{Уравнения возмущений:} Реализована модифицированная иерархия Больцмана \eqref{eq:hier0}--\eqref{eq:hier_ell} с дополнительными членами, пропорциональными $\beta\gamma$, и эффективной оптической толщиной $\tau_{\mathrm{eff}}$. Несохранение энергии-импульса из-за $\dot{G}_{\mathrm{eff}}$ включено в соответствии с формализмом \cite{Uzan2011}.
			\item \textbf{Шкала затухания:} Интеграл силковского затухания пересчитывается с использованием модифицированной скорости звука и вязкости.
			\item \textbf{Вращательные B-моды:} Отдельный модуль вычисляет спектр мощности B-мод, порождаемых глобальным вращением, и добавляет его к стандартным вкладам тензорных мод и линзирования.
		\end{enumerate}
		Новые параметры: $\varepsilon_0$, $\eta$ и амплитуда вихрей $A_v \propto \omega^2$. Для MCMC-анализа мы фиксируем $\beta\gamma = 0,20$ (из Приложения P) и считаем $\varepsilon_0$ и $\eta$ свободными. Космологические параметры ($\Omega_b h^2$, $\Omega_c h^2$, $H_0$, $\tau_{\mathrm{reio}}$, $n_s$, $A_s$) варьируются как обычно. Мы также исследуем корреляцию между новыми параметрами и $n_s$, $A_s$, изучая апостериорные распределения.
		
		\subsection{Вырожденность параметров и систематические эффекты}
		
		Апостериорные распределения показывают слабую положительную корреляцию между $\varepsilon_0$ и $n_s$, а также между $\eta$ и $A_s$. Это ожидаемо, поскольку изменение темпа расширения или шкалы затухания может быть частично скомпенсировано настройкой первичного спектра. Маргинализованные ограничения на стандартные параметры $\Lambda$CDM остаются по существу неизменными, что показывает, что расширение YUCT не вносит сильных напряжений.
		
		Систематические эффекты, которые могут имитировать высоко-$\ell$ избыток, включают остаточное загрязнение точечными источниками, неточное моделирование диаграммы направленности и нелинейные вклады от кинетического эффекта Сюняева--Зельдовича. Тщательный повторный анализ данных Planck с учётом этих эффектов необходим для подтверждения реальности сигнала. Тем не менее, модель YUCT предлагает физически мотивированное (хотя и нестандартное) объяснение, которое можно проверить будущими данными.
		
		\subsection{Статистическая значимость и ценность умеренного улучшения}
		
		Мы отмечаем, что статистическое улучшение по $\chi^2$ является умеренным ($\Delta\chi^2 = -7,6$ для двух дополнительных параметров, $\Delta\mathrm{BIC} \approx -2,2$). Изолированно это не является доказательством новой физики. Однако этот анализ не следует рассматривать как самостоятельное обнаружение. Скорее, он показывает, что \emph{необходимые} поправки, требуемые YUCT (температурная зависимость $G_{\mathrm{eff}}$ и существование координационной вязкости), не портят превосходное согласие $\Lambda$CDM с данными. Более того, они незначительно улучшают его именно в той области, где наблюдался небольшой избыток. Истинная сила YUCT заключается в \emph{одновременном} согласовании этих CMB-данных с независимыми ограничениями от красных смещений белых карликов (Приложение P), системы Земля-Луна (Приложение P) и зависящих от массы отклонений в гравитационно-волновом сигнале затухания (Приложение G). Когда все эти наборы данных рассматриваются совместно, случай для универсального показателя $\beta=0,67$ становится убедительным.
		
		\subsection{Предсказания для будущих экспериментов}
		
		CMB‑S4 с его ожидаемой чувствительностью порядка шума $1\,\mu$K-угл. мин. сможет обнаружить вращательные B-моды с уровнем значимости $>5\sigma$, если параметры YUCT близки к наилучшим значениям. LiteBIRD предоставит критическую перекрёстную проверку на больших угловых масштабах. Кроме того, специфическая картина негауссовости ($f_{\mathrm{NL}}^{\mathrm{local}} \approx 1,1$) находится в пределах досягаемости будущих обзоров крупномасштабной структуры.
		
		\section{Обсуждение и заключение}
		
		Мы представили первую полную обработку мелкомасштабных анизотропий CMB в рамках объединяющей теории координации Якушева. Ключевые результаты:
		\begin{enumerate}
			\item \textbf{Модифицированные уравнения Больцмана}, включающие зависящую от температуры гравитацию и координационную вязкость, выведенные из YUCT-межсекторных связей ($\kappa_{01}$).
			\item \textbf{Связь с алгебраической петлёй YUCT}, показывающая, что новые параметры $\varepsilon_0$ и $\eta$ не произвольны, а связаны с универсальными константами $\beta=2/3$, $\kappa_c=1/3$, $S_{\mathrm{odd}}$ и $S_{\mathrm{even}}$. Это уменьшает эффективное число свободных параметров и усиливает физическую мотивацию.
			\item \textbf{Улучшенное согласие с высоко-$\ell$ данными Planck} с $\Delta\chi^2 = -7,6$ для двух дополнительных параметров. Избыток мощности при $\ell \sim 2000$--$2500$ естественным образом объясняется уменьшенным силковским затуханием из-за координационной вязкости.
			\item \textbf{Уникальное предсказание вращательных B-мод} с характерной формой $\ell^{-2}$, проверяемое в будущих CMB-экспериментах, а также специфическая картина первичной негауссовости.
		\end{enumerate}
		
		Мы обсудили ограничения текущего анализа: статистическое улучшение при изолированном рассмотрении является умеренным, и существуют вырожденности со стандартными параметрами. Систематические эффекты в данных Planck также могут вносить вклад в высоко-$\ell$ остатки. Тем не менее, расширение YUCT обеспечивает когерентную физическую картину, связывающую раннюю Вселенную с лабораторной физикой через универсальный показатель $\beta = 0,67$ и алгебраическую петлю. Истинная сила этого подхода заключается в его способности делать множественные коррелированные предсказания для различных наборов данных (CMB, белые карлики, гравитационные волны), все из которых вытекают из одного набора универсальных констант.
		
		Будущие работы должны включать более детальный расчёт координационной вязкости из полного 19-мерного лагранжиана, а также совместный MCMC-анализ, объединяющий данные CMB, крупномасштабной структуры и астрофизики. Разработанная здесь методология может быть также распространена на наблюдаемые крупномасштабной структуры, обеспечивая всестороннюю проверку координационной парадигмы на всех космологических масштабах.
		
		\section*{Благодарности}
		
		Автор благодарит разработчиков кодов CAMB и MontePython за предоставление их кодов в открытом доступе.
		
		\section*{Доступность данных}
		
		Модифицированный код CAMB и MCMC-цепочки доступны в репозитории YPSDC: \url{https://github.com/Alexey-Yakushev-YUCT/YPSDC}.
		
		\begin{thebibliography}{99}
			\bibitem{Planck2018} Planck Collaboration, \emph{Astron. Astrophys.} \textbf{641}, A6 (2020).
			\bibitem{Ma1995} C.-P. Ma, E. Bertschinger, \emph{Astrophys. J.} \textbf{455}, 7 (1995).
			\bibitem{Silk1968} J. Silk, \emph{Astrophys. J.} \textbf{151}, 459 (1968).
			\bibitem{Dodelson2003} S. Dodelson, \emph{Modern Cosmology}, Academic Press (2003).
			\bibitem{Hu1995} W. Hu, N. Sugiyama, \emph{Astrophys. J.} \textbf{444}, 489 (1995).
			\bibitem{Cooray2005} A. Cooray, D. Baumann, \emph{Phys. Rev. D} \textbf{71}, 063002 (2005).
			\bibitem{Lewis2000} A. Lewis, A. Challinor, A. Lasenby, \emph{Astrophys. J.} \textbf{538}, 473 (2000).
			\bibitem{Audren2013} B. Audren et al., \emph{J. Cosmol. Astropart. Phys.} \textbf{1302}, 001 (2013).
			\bibitem{Uzan2011} J.-P. Uzan, \emph{Living Rev. Relativ.} \textbf{14}, 2 (2011).
			\bibitem{BICEP2021} BICEP/Keck Collaboration, \emph{Phys. Rev. Lett.} \textbf{127}, 151301 (2021).
			\bibitem{AppendixP} A.V. Yakushev, ``Приложение P: Температурная зависимость гравитации,'' в \emph{YUCT}, Zenodo (2026).
			\bibitem{AppendixM} A.V. Yakushev, ``Приложение M: Космология YUCT — Инфляция как координационный фазовый переход,'' в \emph{YUCT}, Zenodo (2026).
			\bibitem{AppendixΩ} A.V. Yakushev, ``Приложение $\Omega$: Глобальное вращение Вселенной и её новый минимальный возраст,'' в \emph{YUCT}, Zenodo (2026).
			\bibitem{AppendixA} A.V. Yakushev, ``Приложение A: Практическое руководство по использованию YUCT V35.0,'' в \emph{YUCT}, Zenodo (2026).
			\bibitem{AppendixY} A.V. Yakushev, ``Приложение Y: Математические основы YUCT,'' в \emph{YUCT}, Zenodo (2026).
			\bibitem{AppendixLambda} A.V. Yakushev, ``Приложение $\Lambda$: Решение проблем иерархии и космологической постоянной в YUCT,'' в \emph{YUCT}, Zenodo (2026).
			\bibitem{AppendixFerra} A.V. Yakushev, ``Приложение Ferra1.2: Универсальные координационные константы для упорядоченной и неупорядоченной фаз,'' в \emph{YUCT}, Zenodo (2026).
		\end{thebibliography}
		
	\end{multicols}
\end{document}